\documentclass[11pt]{article}
\usepackage[margin=1in]{geometry}
\usepackage[T1]{fontenc}
\usepackage[utf8]{inputenc}
\usepackage{lmodern}
\usepackage{amsmath,amssymb,amsthm,mathtools}
\usepackage{enumitem}
\usepackage{booktabs}
\usepackage{hyperref}

\title{Harmonic $t$-ary constructions for hypergraphs with no large exact-one partition}
\author{GPT-5.4 Pro}
\date{}

\newtheorem{theorem}{Theorem}[section]
\newtheorem{proposition}[theorem]{Proposition}
\newtheorem{lemma}[theorem]{Lemma}
\newtheorem{corollary}[theorem]{Corollary}
\newtheorem{definition}[theorem]{Definition}
\newtheorem{remark}[theorem]{Remark}

\DeclareMathOperator{\lcm}{lcm}

\begin{document}
\maketitle

\begin{abstract}
For $n \in \mathbb{N}$, let $H(n)$ denote the maximum number of vertices in a hypergraph with no isolated vertices and with no partition of size greater than $n$ in the sense of Brian and Larson. We give a self-contained lower-bound construction showing that for every fixed integer $t \ge 2$,
\[
H(n) \ge \frac{H_t - 1}{\ln t}\, n \ln n - O_t(n),
\]
where $H_t = \sum_{j=1}^t \frac1j$ is the $t$th harmonic number. The Brian--Larson recurrence is the case $t=2$. Letting $t \to \infty$ after $n \to \infty$ yields
\[
H(n) \ge (1-o(1)) n \ln n.
\]
We also formulate the finite witness problem as an exact integer optimization over ``outer frames,'' which leads to a practical dynamic-programming generator. In particular, this approach produces certified witnesses of sizes $(44,49,52,57,60,65,69,73,75,82,85)$ for $n=15,16,\dots,25$.
\end{abstract}

\section{The hypergraph extremal problem}

A hypergraph is a pair $(V,E)$ with $V$ a finite vertex set and $E$ a finite family of subsets of $V$. Following Brian and Larson, one says that $(V,E)$ has a partition of size $m$ if there exist $D \subseteq V$ and $P \subseteq E$ such that $|D|=m$ and every vertex of $D$ lies in exactly one edge of $P$. Let
\[
H(n) := \max\bigl\{ |V| : (V,E) \text{ has no isolated vertices and no partition of size } > n \bigr\}.
\]

Brian and Larson proved the two-sided estimate
\[
\frac12 n \log_2 n - O(n) < H(n) < n \ln n + \gamma n + \frac12,
\]
and their lower bound is obtained from the binary recurrence
\[
k_1 = 1, \qquad k_n = \Bigl\lfloor \frac n2 \Bigr\rfloor + k_{\lfloor n/2 \rfloor} + k_{\lceil n/2 \rceil}.
\]
The purpose of this note is to isolate the hypergraph problem and to prove an asymptotically stronger lower bound.

The main result is the following.

\begin{theorem}[1.1]
For every fixed integer $t \ge 2$, if $H_t = 1 + \frac12 + \cdots + \frac1t$ and
\[
c_t := \frac{H_t - 1}{\ln t},
\]
then there is a constant $D_t$ such that
\[
H(n) \ge c_t n \ln n - D_t n \qquad \text{for all } n \in \mathbb{N}.
\]
Consequently,
\[
H(n) \ge (1-o(1)) n \ln n.
\]
\end{theorem}

The proof is entirely constructive. The key ingredient is a multiway composition lemma in which the ``outer layer'' is a multiset of supports on the block indices, chosen so that the vertices counted exactly once in the outer layer can be paid for by the unused partition budget of the blocks in which no edge was selected.

\section{Support multisets and exact-one counts}

Fix a hypergraph $H=(V,E)$ and enumerate the edges as $E = \{e_1,\dots,e_m\}$. Each vertex $v \in V$ has an incidence support
\[
\operatorname{supp}(v) := \{ i \in [m] : v \in e_i \} \subseteq [m].
\]
Since isolated vertices are excluded, every support is nonempty. Thus $H$ may be identified with a finite multiset $\mathcal S$ of nonempty subsets of $[m]$.

For $P \subseteq [m]$, define the exact-one count
\[
\operatorname{ex1}_{\mathcal S}(P) := \bigl|\{ S \in \mathcal S : |S \cap P| = 1\}\bigr|,
\]
with multiplicities counted. This is exactly the number of vertices lying in precisely one selected edge.

\begin{proposition}[2.1]
A hypergraph $H$ has no partition of size greater than $n$ if and only if
\[
\operatorname{ex1}_{\mathcal S}(P) \le n \qquad \text{for every } P \subseteq [m],
\]
where $\mathcal S$ is the multiset of supports of the vertices of $H$.
\end{proposition}

\begin{proof}
If $D \subseteq V$ and $P \subseteq E$ witness a partition of size $m$, then every vertex of $D$ is counted by $\operatorname{ex1}_{\mathcal S}(P)$, so $\operatorname{ex1}_{\mathcal S}(P) \ge m$. Conversely, if $\operatorname{ex1}_{\mathcal S}(P)=m$, then one may take $D$ to be the set of all vertices counted exactly once by $P$. Thus the largest partition size is exactly $\max_P \operatorname{ex1}_{\mathcal S}(P)$.
\end{proof}

Henceforth we work in the support-set model. A witness for $H(n)$ is simply a multiset $\mathcal S$ of nonempty subsets of some edge set $[m]$ such that $\operatorname{ex1}_{\mathcal S}(P) \le n$ for every $P \subseteq [m]$.

\section{Outer frames and a composition lemma}

Let $t \ge 2$. Suppose we have $t$ disjoint edge blocks
\[
[m_1], [m_2], \dots, [m_t]
\]
(which later will support recursive witnesses), and let $\mathcal F$ be a multiset of nonempty subsets of $[t]$. Every $A \in \mathcal F$ determines an outer support
\[
U_A := \bigcup_{i \in A} [m_i],
\]
that is, a vertex incident with every edge in each block indexed by $A$.

To control exact-one counts, the correct constraint is not on all subsets of $[t]$, but on the patterns of blocks containing $0$, $1$, or at least $2$ chosen edges.

\begin{definition}[3.1]
Let $\mathbf n = (n_1,\dots,n_t) \in \mathbb N^t$. A multiset $\mathcal F$ of nonempty subsets of $[t]$ is $\mathbf n$-admissible if for every pair of disjoint subsets $Z,I \subseteq [t]$ with $I \neq \varnothing$,
\[
N_{\mathcal F}(Z,I) := \bigl|\{ A \in \mathcal F : A \subseteq Z \cup I \text{ and } |A \cap I|=1\}\bigr| \le \sum_{i \in Z} n_i.
\]
Multiplicities are counted in $N_{\mathcal F}(Z,I)$.
\end{definition}

The interpretation is simple: $I$ records the blocks in which exactly one edge was selected, $Z$ records the blocks in which no edge was selected, and the right-hand side is the unused partition budget of the zero-selected blocks.

\begin{lemma}[3.2, Composition lemma]
For $1 \le i \le t$, let $\mathcal S_i$ be a support multiset on a block of edges $E_i$ such that
\[
\operatorname{ex1}_{\mathcal S_i}(P) \le n_i \qquad \text{for every } P \subseteq E_i.
\]
Assume that the edge sets $E_1,\dots,E_t$ are pairwise disjoint, and let $\mathcal F$ be a $\mathbf n$-admissible multiset on $[t]$. Form the combined support multiset
\[
\mathcal S := \mathcal S_1 \sqcup \cdots \sqcup \mathcal S_t \sqcup \{ U_A : A \in \mathcal F \},
\]
where $U_A = \bigcup_{i \in A} E_i$. Then
\[
\operatorname{ex1}_{\mathcal S}(P) \le n_1 + \cdots + n_t \qquad \text{for every } P \subseteq E_1 \sqcup \cdots \sqcup E_t.
\]
In particular,
\[
H(n_1 + \cdots + n_t) \ge |\mathcal S_1| + \cdots + |\mathcal S_t| + |\mathcal F|.
\]
\end{lemma}

\begin{proof}
Fix $P \subseteq E_1 \sqcup \cdots \sqcup E_t$, and write $P_i = P \cap E_i$. Partition $[t]$ into
\[
Z := \{ i : |P_i| = 0\}, \qquad I := \{ i : |P_i| = 1\}, \qquad M := \{ i : |P_i| \ge 2\}.
\]
The contribution from the inner block $\mathcal S_i$ is at most $n_i$ whenever $P_i \neq \varnothing$, and is $0$ when $P_i = \varnothing$. Hence the inner blocks contribute at most
\[
\sum_{i \in I \cup M} n_i.
\]
Now consider an outer support $U_A$ coming from some $A \in \mathcal F$. This support is counted exactly once by $P$ if and only if the total number of selected edges meeting it is $1$, i.e.
\[
\sum_{i \in A} |P_i| = 1.
\]
Equivalently, $A \subseteq Z \cup I$ and $|A \cap I| = 1$. Thus the outer supports contribute exactly $N_{\mathcal F}(Z,I)$, which is at most $\sum_{i \in Z} n_i$ by admissibility. Therefore
\[
\operatorname{ex1}_{\mathcal S}(P) \le \sum_{i \in I \cup M} n_i + \sum_{i \in Z} n_i = \sum_{i=1}^t n_i.
\]
This proves the first claim. The final inequality follows by choosing each $\mathcal S_i$ to be a witness for $H(n_i)$.
\end{proof}

\begin{remark}[3.3]
The Brian--Larson binary construction is exactly the case $t=2$ with $\mathcal F$ consisting of some number of copies of $\{1,2\}$. Definition~3.1 then says that the number of such copies may be at most $\min(n_1,n_2)$.
\end{remark}

\section{Harmonic $t$-ary frames}

The previous lemma reduces the problem to the construction of good admissible outer frames. The following family is the symmetric choice that gives the best asymptotic coefficient we presently know how to justify.

Fix $t \ge 2$ and let
\[
L_t := \operatorname{lcm}\left\{ \binom{t-1}{j} : 1 \le j \le t-1 \right\}.
\]
If $q$ is a positive multiple of $L_t$, define the multiset $\mathcal F_t(q)$ on $[t]$ by putting
\[
\frac{q}{\binom{t-1}{s-1}}
\]
copies of every $s$-subset of $[t]$, for each $s=2,3,\dots,t$.

The proof of admissibility rests on a simple binomial identity.

\begin{lemma}[4.1]
For integers $0 \le z \le N$,
\[
\sum_{j=0}^z \binom{z}{j} \binom{N}{j}^{-1} = \frac{N+1}{N-z+1}.
\]
Consequently,
\[
\sum_{j=1}^z \binom{z}{j} \binom{N}{j}^{-1} = \frac{z}{N-z+1}.
\]
\end{lemma}

\begin{proof}
Using
\[
\binom{z}{j} \binom{N}{z}^{-1} = \binom{N}{j}^{-1} \binom{N-j}{z-j},
\]
we obtain
\[
\sum_{j=0}^z \binom{z}{j} \binom{N}{j}^{-1} = \binom{N}{z}^{-1} \sum_{j=0}^z \binom{N-j}{z-j}.
\]
Now set $r = z-j$. Then the sum becomes
\[
\sum_{r=0}^z \binom{N-z+r}{r} = \binom{N+1}{z}
\]
by the hockey-stick identity. Dividing by $\binom{N}{z}$ gives
\[
\frac{\binom{N+1}{z}}{\binom{N}{z}} = \frac{N+1}{N-z+1}.
\]
Subtracting the $j=0$ term gives the second formula.
\end{proof}

\begin{proposition}[4.2]
Let $t \ge 2$ and let $q$ be a positive multiple of $L_t$. Then $\mathcal F_t(q)$ is $(q,\dots,q)$-admissible. Moreover,
\[
|\mathcal F_t(q)| = tq(H_t - 1).
\]
\end{proposition}

\begin{proof}
Fix disjoint $Z,I \subseteq [t]$ with $I \neq \varnothing$, and write $z = |Z|$ and $i = |I|$. For a set $A \in \mathcal F_t(q)$ to contribute to $N_{\mathcal F_t(q)}(Z,I)$, it must contain exactly one point of $I$ and any number $j \ge 1$ of points of $Z$. Thus
\[
N_{\mathcal F_t(q)}(Z,I) = q i \sum_{j=1}^z \binom{z}{j} \binom{t-1}{j}^{-1}.
\]
Applying Lemma~4.1 with $N=t-1$ gives
\[
N_{\mathcal F_t(q)}(Z,I) = q i \frac{z}{t-z}.
\]
Since $Z$ and $I$ are disjoint subsets of $[t]$, we have $i \le t-z$, and therefore
\[
N_{\mathcal F_t(q)}(Z,I) \le qz = \sum_{u \in Z} q.
\]
This is exactly the admissibility condition.

For the size, we compute
\[
|\mathcal F_t(q)| = \sum_{s=2}^t \binom{t}{s} \frac{q}{\binom{t-1}{s-1}} = q \sum_{s=2}^t \frac{t}{s} = tq\left(\frac12 + \frac13 + \cdots + \frac1t\right) = tq(H_t - 1).
\]
\end{proof}

Combining Proposition~4.2 with Lemma~3.2 yields the basic multiway recurrence.

\begin{corollary}[4.3]
Fix $t \ge 2$. Let $n_1,\dots,n_t \in \mathbb N$, and let $q$ be any multiple of $L_t$ with $0 \le q \le \min_i n_i$. Then
\[
H(n_1 + \cdots + n_t) \ge H(n_1) + \cdots + H(n_t) + tq(H_t - 1).
\]
\end{corollary}

\begin{remark}[4.4]
For $t=2$ we have $L_2=1$ and $H_2-1=\frac12$, so Corollary~4.3 becomes
\[
H(n_1+n_2) \ge H(n_1) + H(n_2) + \min(n_1,n_2),
\]
which is exactly the Brian--Larson binary recurrence. For $t=4$, $L_4=3$, and $q=3$ gives one copy of each $2$-subset, one copy of each $3$-subset, and three copies of $[4]$, i.e. the $13$-vertex outer gadget discussed earlier in this line of work.
\end{remark}

\section{Asymptotics for fixed arity}

We now prove the fixed-$t$ lower bound.

\begin{theorem}[5.1]
Fix $t \ge 2$, and let
\[
c_t := \frac{H_t - 1}{\ln t}.
\]
Then there exists a constant $D_t > 0$ such that
\[
H(n) \ge c_t n \ln n - D_t n \qquad \text{for every } n \in \mathbb N.
\]
\end{theorem}

\begin{proof}
Fix $t$. Set
\[
K_t := (H_t - 1)(tL_t - 1), \qquad A_t := \frac{K_t}{t-1}.
\]
We claim that there exists $D_t$ such that
\begin{equation}
H(n) \ge c_t n \ln n - D_t n + A_t \qquad \text{for all } n \in \mathbb N. \tag{1}
\end{equation}
Since $H(n) \ge n$ (take $n$ singleton edges and one vertex on each edge), it is enough to choose $D_t$ so large that
\[
n \ge c_t n \ln n - D_t n + A_t \qquad \text{for every } 1 \le n < tL_t.
\]
This is possible because only finitely many $n$ are involved.

We now prove \eqref{1} by induction on $n$. The base range $1 \le n < tL_t$ has already been absorbed into the choice of $D_t$. Assume next that $n \ge tL_t$, and write
\[
n = tm + r, \qquad 0 \le r < t.
\]
Define
\[
q := L_t \left\lfloor \frac{m}{L_t} \right\rfloor.
\]
Then $q$ is a multiple of $L_t$ and $q \le m \le m+1$. Applying Corollary~4.3 to $r$ blocks of size $m+1$ and $t-r$ blocks of size $m$ gives
\begin{equation}
H(n) \ge rH(m+1) + (t-r)H(m) + tq(H_t - 1). \tag{2}
\end{equation}
Because $q \ge m - (L_t - 1)$, we have
\[
tq = t\bigl(m - (m-q)\bigr) = n-r-t(m-q) \ge n-(t-1)-t(L_t-1) = n-(tL_t-1).
\]
Hence
\begin{equation}
tq(H_t - 1) \ge (H_t - 1)n - K_t. \tag{3}
\end{equation}
Now apply the induction hypothesis to the terms on the right side of \eqref{2}. Since $r(m+1) + (t-r)m = n$, we obtain
\[
H(n) \ge c_t\bigl(r(m+1)\ln(m+1) + (t-r)m\ln m\bigr) - D_t n + tA_t + tq(H_t - 1).
\]
The function $x \mapsto x\ln x$ is convex on $(0,\infty)$, so by Jensen's inequality,
\[
r(m+1)\ln(m+1) + (t-r)m\ln m \ge n\ln\left(\frac{n}{t}\right).
\]
Combining this with \eqref{3} gives
\begin{align*}
H(n)
&\ge c_t n \ln\left(\frac nt\right) - D_t n + tA_t + (H_t - 1)n - K_t \\
&= c_t n \ln n - c_t n \ln t - D_t n + (H_t - 1)n + tA_t - K_t.
\end{align*}
Since $c_t \ln t = H_t - 1$ and $(t-1)A_t = K_t$, the last line simplifies to
\[
H(n) \ge c_t n \ln n - D_t n + A_t.
\]
This is exactly \eqref{1}, completing the induction.

Finally, dropping the positive constant term $A_t$ yields the stated bound.
\end{proof}

\begin{corollary}[5.2]
One has
\[
H(n) \ge (1-o(1)) n \ln n.
\]
\end{corollary}

\begin{proof}
Let $\varepsilon > 0$. Because $H_t = \ln t + \gamma + o(1)$ as $t \to \infty$, the coefficients
\[
c_t = \frac{H_t - 1}{\ln t}
\]
converge to $1$. Choose $t$ so large that $c_t > 1 - \varepsilon/2$. Then Theorem~5.1 gives a constant $D_t$ such that
\[
H(n) \ge c_t n \ln n - D_t n
\]
for all $n$. For sufficiently large $n$, the linear error term satisfies $D_t n \le \frac{\varepsilon}{2} n \ln n$. Hence
\[
H(n) \ge (1-\varepsilon)n \ln n
\]
for all sufficiently large $n$. Since $\varepsilon > 0$ was arbitrary, the result follows.
\end{proof}

\begin{remark}[5.3]
The coefficient in Theorem~5.1 may also be written in base-$2$ notation as
\[
\frac{H_t - 1}{\log_2 t}\, n \log_2 n.
\]
Thus $t=2$ gives the Brian--Larson leading coefficient $\frac12$ in base $2$, while $t=4$ gives $\frac{13}{24}$ in base $2$.
\end{remark}

\section{Practical generation of witnesses}

The asymptotic proof above produces an explicit witness generator: recursively split the edge budget into $t$ nearly equal parts, use the harmonic frame $\mathcal F_t(q)$ as the outer layer, and stop at a small library of certified seeds. In practice one can do better at moderate $n$ by solving exactly for the best outer frame over a fixed profile.

\subsection{Exact optimization over outer frames}

Fix positive integers $n_1,\dots,n_t$. Introduce one nonnegative integer variable $x_A$ for each subset $A \subseteq [t]$ with $|A| \ge 2$; $x_A$ represents the number of outer vertices supported on the union of the blocks indexed by $A$. By Definition~3.1, the admissibility constraints are
\begin{equation}
\sum_{\substack{A \subseteq Z \cup I \\ |A \cap I| = 1}} x_A \le \sum_{i \in Z} n_i
\qquad \text{for every disjoint } Z,I \subseteq [t] \text{ with } I \neq \varnothing. \tag{4}
\end{equation}
Every feasible integer solution of \eqref{4} defines a valid outer frame, and therefore Lemma~3.2 gives the certified lower bound
\[
H(n_1 + \cdots + n_t) \ge H(n_1) + \cdots + H(n_t) + \sum_{|A| \ge 2} x_A.
\]
The harmonic frame of Section~4 is exactly the symmetric feasible point
\[
x_A = \frac{q}{\binom{t-1}{|A|-1}}, \qquad |A| \ge 2,
\]
with $q$ a multiple of $L_t$ and $q \le \min_i n_i$.

For fixed $t$, the optimization problem is very small: it has $2^t - t - 1$ variables and $3^t - 2^t$ inequalities. In particular, exact search is practical for $t \le 7$ or $8$, and that is already enough to improve substantially on the binary lower bound in the range of interest here.

\subsection{Dynamic programming and reconstruction}

A practical witness generator works recursively as follows.

\begin{enumerate}[leftmargin=2em]
\item Store a library of certified seed witnesses for small $n$.
\item For each target $n$, enumerate a menu of block profiles $n = n_1 + \cdots + n_t$ (balanced splits are a very good default; exact outer-frame search over a few small arities improves the small values).
\item For each profile, solve \eqref{4} to obtain a certified outer bonus, add the previously computed values of the subproblems, and retain the best option.
\item Store back-pointers consisting of the chosen profile and the chosen outer frame. Recursive expansion of these back-pointers produces an explicit support multiset on an edge set of size $n$.
\item Translate the support multiset into the requested edge-list format by placing vertex $v$ in precisely those edges whose labels belong to its support.
\end{enumerate}

The recursive proof certificate is built into the data structure: each back-pointer is justified by Lemma~3.2, so a full witness can be verified recursively without enumerating all $2^n$ subsets of edges. For $n \le 25$, however, brute-force verification is also feasible and provides an additional independent check.

\subsection{Certified witnesses for $15 \le n \le 25$}

A machine-assisted search over small outer frames (with arity at most $7$) produces the following certified witness sizes. They are not claimed to be globally optimal, but each line is accompanied by an explicit decomposition and hence by a complete recursive correctness proof.

\begin{center}
\begin{tabular}{@{}cccc@{}}
\toprule
$n$ & Brian--Larson $k_n$ & certified $|V|$ & gain \\
\midrule
15 & 43 & 44 & 1 \\
16 & 48 & 49 & 1 \\
17 & 50 & 52 & 2 \\
18 & 53 & 57 & 4 \\
19 & 56 & 60 & 4 \\
20 & 60 & 65 & 5 \\
21 & 63 & 69 & 6 \\
22 & 67 & 73 & 6 \\
23 & 71 & 75 & 4 \\
24 & 76 & 82 & 6 \\
25 & 79 & 85 & 6 \\
\bottomrule
\end{tabular}
\end{center}

For example, the $n=20$ witness is obtained from a five-way profile $(4,4,4,4,4)$ and an exact outer frame on those five blocks, while the $n=18$ witness is obtained from the profile $(2,2,2,6,6)$. The explicit back-pointers needed to reconstruct these witnesses are compact, and a companion Python script can emit the final hypergraph as a string in the edge-list format
\[
\{1,2,3\},\ \{2,4\},\ \{3,4,5\},\dots
\]
required by the original problem statement.

\section{Concluding remarks}

The support-set viewpoint makes the lower-bound problem almost entirely transparent: one is trying to build a large multiset of nonempty subsets of an edge set subject to the exact-one inequalities
\[
\operatorname{ex1}_{\mathcal S}(P) \le n \qquad (P \subseteq [m]).
\]
The outer-frame formulation isolates the recursive step from the internal structure of the blocks, and the harmonic frame $\mathcal F_t(q)$ gives a particularly clean family of admissible solutions. Theorem~5.1 shows that each fixed arity $t$ contributes the coefficient $(H_t - 1)/\ln t$, and Corollary~5.2 shows that the leading constant can be driven all the way up to $1$.

Thus the hypergraph problem admits explicit witnesses on
\[
(1-o(1)) n \ln n
\]
vertices. This matches the leading constant in the Brian--Larson upper bound and places the remaining difficulty in the lower-order terms and in the finite optimization of the outer frames.

\begin{thebibliography}{9}
\bibitem{BrianLarson2021}
W.~Brian and P.~B.~Larson,
\emph{Choosing between incompatible ideals},
European Journal of Combinatorics \textbf{96} (2021), 103349.
\end{thebibliography}

\end{document}
